\section{An elementary witness for every index}\label{sec:elementary}
Write $v_q$ for the $q$-adic valuation. For $m\ge1$ put
$S_p(m)=1+p+\cdots+p^{m-1}$, so the mean of the
divisors of $p^{m-1}$ is $S_p(m)/m$.

\begin{lemma}[Upward numerator primes]\label{lem:upward}
Let $q,p$ be primes, $p\not\equiv1\pmod q$, and $b\ge1$.
Every prime divisor of $S_p(q^b)$ is strictly greater than $q$.
\end{lemma}
\begin{proof}
The sum is 1 modulo $p$. It is also 1 modulo $q$: if $p=q$ this is
immediate, and otherwise Fermat's theorem gives
$p^{q^b}\equiv p\pmod q$, with $p-1$ invertible modulo $q$.
If a prime $r$ divides both the sum and $p-1$, reducing the sum modulo
$r$ gives $r\mid q^b$, contradicting the preceding observation.
For any other prime divisor $r$, the order of $p$ modulo $r$ is a
nontrivial divisor of $q^b$. Hence $q\mid r-1$, and $r>q$.
\end{proof}

\begin{theorem}[Explicit elementary upper bound]\label{thm:elementary}
Write $n=\prod_{i=1}^s q_i^{a_i}$ with $q_1<\cdots<q_s$, and let
$\ell_i$ be the $i$-th prime. Define $U_0=1$ and, successively,
\begin{equation}\label{eq:elementary-bound}
 b_i=a_i+v_{q_i}(\sigma(U_{i-1})),\qquad
 U_i=U_{i-1}\ell_i^{q_i^{b_i}-1}.
\end{equation}
Then $\den(\sigma(U_s)/\tau(U_s))=n$, so $a(n)\le U_s$.
For $n=1$ use $U_0=1$.
\end{theorem}
\begin{proof}
Since $\ell_i\le q_i$, the new base prime is not 1 modulo $q_i$ and
has not been used before. Its divisor count is $q_i^{b_i}$, and by
Lemma~\ref{lem:upward} its divisor sum has no prime factor at most $q_i$.
It therefore leaves all smaller denominator valuations unchanged.
The previous divisor count is supported on $q_1,\ldots,q_{i-1}$,
so the new net $q_i$-valuation of the denominator is exactly
$b_i-v_{q_i}(\sigma(U_{i-1}))=a_i$. Later factors do not change this
valuation. The final divisor count has no primes outside those of $n$,
which proves the assertion. This construction needs no theorem about
primes in arithmetic progressions.
\end{proof}

\section{A divisor-preserving normal form}\label{sec:normal}
For primes $q,p$ and $m\ge1$, define the signed local contribution
\[
 c_q(p,m)=v_q(m)-v_q(S_p(m)).
\]
For $K=\prod p_i^{m_i-1}$, the valuation of $q$ in its reduced
denominator is $\max(0,\sum_i c_q(p_i,m_i))$.

\begin{lemma}[Triangular dependence]\label{lem:triangular}
For an odd prime $q$, the contributions are as follows, with
$h=\operatorname{ord}_q(p)$ in the last two cases:
\begin{equation}\label{eq:local-contributions}
 c_q(p,m)=
 \begin{cases}
 v_q(m),&p=q,\\
 0,&p\equiv1\pmod q,\\
 v_q(m),&p\not\equiv0,1\pmod q,\ h\nmid m,\\
 -v_q(p^h-1),&p\not\equiv0,1\pmod q,\ h\mid m,
 \end{cases}
\end{equation}
For $q=2$, $c_2(2,m)=v_2(m)$, while for odd $p$ the contribution is
0 when $m$ is odd and $1-v_2(p+1)$ when $m$ is even.
Deleting factors of primes smaller than $q$ from $m$ cannot decrease
$c_q(p,m)$; deleting factors of primes greater than $q$ cannot change it.
\end{lemma}
\begin{proof}
The cases $p=q$ and $h\nmid m$ have numerator valuation zero.
For $p\equiv1\pmod q$, the odd-prime lifting identity gives
$v_q(S_p(m))=v_q(m)$. If $h\mid m$ and $p\not\equiv0,1\pmod q$,
lifting gives $v_q(S_p(m))=v_q(p^h-1)+v_q(m)$.
The stated 2-adic identity follows from the usual even-exponent
lifting formula, and an odd number of odd summands is odd.
Finally $h\mid q-1$. Deleting a smaller prime can only change
$h\mid m$ to $h\nmid m$, replacing a negative contribution by a
nonnegative one. Deleting a larger prime changes neither this condition
nor $v_q(m)$. The assertions for $q=2$ follow from parity.
\end{proof}

\begin{theorem}[Reduced-factor normal form]\label{thm:normal}
If $\den(\sigma(K)/\tau(K))=n$, there is a divisor
$K'=\prod p_i^{m_i-1}$ of $K$ with the same reduced denominator such that
\begin{enumerate}
\item every prime divisor of every $m_i$ divides $n$;
\item $\gcd(m_i,S_{p_i}(m_i))=1$ for every $i$.
\end{enumerate}
Consequently the least value $a(n)$ itself has both properties.
\end{theorem}
\begin{proof}
Allow $m_i=1$ to retain a base with exponent zero. Process the prime
divisors of the original $\tau(K)=\prod m_i$ in increasing order,
only replacing $m_i$ by divisors of themselves.

At $q=2$, remove all 2-parts from the $m_i$ attached to odd bases.
If $2\mid n$, retain exactly $v_2(n)$ powers of 2 in the $m_i$ for
base 2; the original exponent has enough capacity, since every odd
base contributes nonpositively at 2. If $2\nmid n$, remove its 2-part
as well. All local means are now reduced at 2, and the denominator has
the required 2-valuation. Larger denominator primes have not lost
any signed contribution, by Lemma~\ref{lem:triangular}.

For an odd prime $q$, classify the current factors using
\eqref{eq:local-contributions}. Call the factors in the first or third
case active. The fourth-case contributions are constants $-\alpha_i$,
and the second-case contributions are zero. Remove the $q$-parts of
$m_i$ from every inactive factor; this changes none of those signed
contributions. If $q\nmid n$, also remove all active $q$-parts, leaving
a nonpositive sum and hence denominator valuation zero.
If $q\mid n$, put $a=v_q(n)$. Earlier steps have only increased the
signed $q$-sum, so the active capacities satisfy
\[
 \sum_{i\ \mathrm{active}}v_q(m_i)\ \ge\ a+\sum_{i\ \mathrm{fourth}}\alpha_i.
\]
Retain any allocation of exactly the right-hand side among these
capacities. The resulting signed sum is exactly $a$.
Active numerator sums are coprime to $q$, and the other factors now
have $q$-free $m_i$, so every local mean is reduced at $q$.
Later steps leave this conclusion and the total $q$-contribution
unchanged. Earlier denominator valuations also remain unchanged.

At the end, the divisor count is supported on the primes of $n$, all
local means are reduced, and the total reduced denominator is $n$.
All exponent changes decreased $m_i$, so $K'\mid K$.
Applying this to a least $K$ forces $K'=K$.
\end{proof}

\begin{corollary}\label{cor:square}
The odd part of $a(n)$ is always a square. If $n$ is odd, $a(n)$ is
itself a square. If $2^a\Vert n$ with $a\ge1$, then $2\mid a(n)$ and
\[
 v_2\bigl(v_2(a(n))+1\bigr)=a.
\]
\end{corollary}
\begin{proof}
An odd base with even $m_i$ violates local coprimality at 2. Thus every
odd-base exponent $m_i-1$ is even. If $n$ is odd, all $m_i$ are odd by
the support restriction. In the remaining case only base 2 can supply
the required 2-valuation, and all other local contributions at 2 are zero.
\end{proof}

\section{The exact minimum for every index}\label{sec:all-minima}
The following is a finite arithmetic formula using multiplicative orders,
valuations, weak compositions, and a finite minimum. Its stages depend
only on the factorization of $n$, not on previous sequence values. It is
not asserted to be an elementary expression without finite minimization.

For $n>1$, fix the explicit witness $U=U_s$ of
\eqref{eq:elementary-bound} and set $L=\lfloor\log_2 U\rfloor$.
Take each tuple of distinct primes
\begin{equation}\label{eq:prime-tuples}
 2\le p_1<\cdots<p_t\le U,\qquad 1\le t\le L.
\end{equation}
Start with $m_1=\cdots=m_t=1$. For $q=q_1,\ldots,q_s$ in increasing
order, write $a=v_q(n)$ and use the current $m_i$ to form
\begin{align*}
 G_q={}&\{i:p_i=q\}\ \cup
 \{i:p_i\not\equiv0,1\pmod q,\ \operatorname{ord}_q(p_i)\nmid m_i\},\\
 N_q={}&\{i:p_i\not\equiv0,1\pmod q,\ \operatorname{ord}_q(p_i)\mid m_i\},\\
 B_q={}&a+\sum_{i\in N_q}v_q(p_i^{\operatorname{ord}_q(p_i)}-1).
\end{align*}
For $q=2$, these definitions simply give $G_2=\{i:p_i=2\}$,
$N_2=\varnothing$, and $B_2=v_2(n)$.
For every weak composition
\begin{equation}\label{eq:allocations}
 b_i\ge0\ (i\in G_q),\qquad \sum_{i\in G_q}b_i=B_q,
\end{equation}
make a branch with $m_i\leftarrow m_iq^{b_i}$ on $G_q$, leaving the
other $m_i$ unchanged. If $G_q$ is empty, there is no branch.
Let $\mathcal C(n)$ consist of all products $\prod p_i^{m_i-1}$ obtained
after the last stage, over the tuples \eqref{eq:prime-tuples} and all
the specified weak compositions. Exponents zero may be discarded.

\begin{theorem}[All-index minimum formula]\label{thm:all-minimum}
The set $\mathcal C(n)$ is finite and nonempty, every member has reduced
divisor-mean denominator $n$, and
\begin{equation}\label{eq:all-minimum}
 \boxed{a(1)=1,\qquad a(n)=\min\mathcal C(n)\quad(n>1).}
\end{equation}
Any verified witness may replace $U$ in \eqref{eq:prime-tuples}.
\end{theorem}
\begin{proof}
Before the $q$-stage, $m_i$ has only prime factors smaller than $q$.
Adding $q$-powers changes the active contributions to $b_i$; the
inactive contributions remain zero or the constants occurring in $B_q$.
Equation~\eqref{eq:allocations} therefore makes the signed sum exactly
$a$. At 2 only base 2 receives an even $m_i$, so the same conclusion
holds. All local means remain reduced at the processed prime.
Later stages have larger primes and cannot change any earlier signed
contribution. Since the final divisor count is supported on the primes
of $n$, every output has denominator exactly $n$.

There are finitely many prime tuples, finitely many stages, and finitely
many weak compositions at each stage. Thus the set is finite.
By Theorem~\ref{thm:elementary}, a least value $K=a(n)\le U$ exists.
Its number of distinct prime factors is at most $\lfloor\log_2 U\rfloor$,
and each base prime is at most $U$. By Theorem~\ref{thm:normal}, its
exponent-plus-one values have the required support and local coprimality.
At every stage, local coprimality forces $b_i=0$ on inactive factors,
while its denominator valuation forces exactly \eqref{eq:allocations}.
Its factorization is consequently one of the branches in $\mathcal C(n)$.
This proves nonemptiness, completeness for the least value, and the
formula. Only $a(n)\le U$ was needed, giving the last assertion.
\end{proof}

\subsection{Exact evaluation without scanning successive integers}
The implementation searches increasing tuples of base primes with
$m_i$ supported on the primes of $n$ and checks local coprimality.
For a partial product $V$, let $c_q$ be its current signed contribution
and put
\[
 E=\sum_{q\mid n}\max(0,v_q(n)-c_q)(q-1).
\]
Any completion needs at least $E$ additional prime factors counted
with multiplicity: for every positive $m$ supported on these primes,
\[
 m-1\ge\sum_{q\mid n}v_q(m)(q-1),
\]
by expanding the product of its prime factors, and additional numerator
valuations can only increase the required divisor-count contributions.
If the partial denominator is not already $n$, at least one additional
factor is required, even if $E=0$. When the next base prime is $p$, the
completion is therefore at least $Vp^{\max(1,E)}$. This supplies a
rigorous stopping bound for each prime loop. Exponents obey
$m_i\le1+\lfloor\log_{p_i}((U-1)/V)\rfloor$ when seeking an improvement
to a known witness $U$. Every cutoff uses exact integer comparisons.

For smaller starting bounds, the implementation may split a $q$-budget
among several unused primes not congruent to 1 modulo $q$, using
Lemma~\ref{lem:upward}. Such primes are elementary to obtain for odd
$q$: if a finite list exhausted them, a prime divisor not congruent to 1
modulo $q$ of $q$ times their product minus 1 would be missing from the
list. The case $q=2$ uses only base 2. These choices affect speed, not
the completeness theorem; every starting bound is independently checked.

For example, $n=30$ and bases $(2,3,5)$ admit the successive budgets
$B_2=1$, $B_3=2$, and $B_5=1$. Choosing allocations $(1,0,0)$,
$(0,1,1)$, and $(1,0,0)$ in base order, with zeros on inactive bases,
gives final $m=(10,3,3)$, hence
\[
 a(30)=2^9 3^2 5^2=115200.
\]
The exact evaluator excludes all smaller candidates using the proved
normal form and bounds. In particular, the minimum cannot be obtained
by requiring $\tau(K)=n$: here $\tau(K)=90$ and cross-factor cancellation
is essential. Further examples are
$a(40)=6480000$, $a(42)=1843200$, $a(45)=3600$, and $a(60)=4706440200$.
